\begin{afterword}
\summaryhead

The post argues that every organ, the brain's circuits included, exists because it helped ancestors reproduce. Anger exists because angry ancestors had more children, and there is no other way it could have got there. Many people hear this as a claim that they are secretly trying to have children when they get angry. That is a confusion.

The cause of an adaptation, its shape and its consequence are separate things. A toaster knows nothing of its designer's purpose. A human is harder to think about, because a human is also a mind with purposes of its own, as the word ``blue'' printed in another color is harder to name than ``wind.'' Cognitive causes are made of neurons; evolutionary causes are made of ancestors. No brain represented inclusive genetic fitness before the twentieth century. One must keep apart an episode of anger, the circuitry, the genes that build it and the ancestral history that explains them.

A story of two strangers in a bar, each drawn by cues that once signalled fertility or status, who use a condom and part, shows adaptations executing with no thought of reproduction. The post ends by saying that evolutionary psychology demands that all these distinctions be kept without a single slip.

\responsehead

The post's central distinction is correct, and it is old: it is Mayr's ultimate and proximate causation, discussed under ``Adaptation-Executers, not Fitness-Maximizers.'' The post explains it at length, through a toaster, the Stroop test, a hand, and a story in a bar. Its most exact statement is a list near the end: the episode of anger, the circuitry, the genes, and the ancestral history.

One claim goes further than the rest. ``Anger exists in Homo sapiens because angry ancestors had more kids. There's no other way it could have gotten there.'' A trait can also be a by-product of selection for something else, and the author allows as much in a later post, which calls the pain of sunk costs possibly ``an unfortunate spandrel.'' Later research supports the view that anger is an adaptation, but this post gives no evidence for it and rules out the alternative by assertion.

The mistake the post corrects is attributed to ``all too many people,'' and no one who makes it is named or cited.

In short: a correct and long-established distinction, explained at length, with one overstatement about how anger arose and an opponent who is not identified.
\end{afterword}
